\subsection{The exact additional condition across a nilpotent thickening}

\begin{lemma}[Serre--Tate across a finite nilpotent thickening of a complete base]
\label{lem:C-RD12-complete}
Let $R$ be a noetherian $p$-adically complete ring with $p$ in its
Jacobson radical, let $J\subset R$ be nilpotent, and put $R_0=R/J$.
Fix an abelian scheme $A_0/R_0$, a BT group $M/R$, and an
isomorphism $A_0[p^\infty]\simeq M_{R_0}$. There is an abelian scheme
$A/R$ with these prescribed reductions, unique up to a unique
isomorphism respecting the prescribed reduction and BT comparison.
No polarization is required. This assertion concerns a finite
nilpotent thickening of an already algebraic abelian scheme, not
algebraization of an arbitrary infinite formal abelian deformation.
\end{lemma}
\begin{proof}
First suppose $J^2=0$. Actual abelian liftings of $A_0$ are locally
nonempty and form a torsor under
\[
 \mathcal F=\operatorname{Lie}(A_0)\otimes
             \operatorname{Lie}(A_0^\vee)\otimes J.
\]
Here and below a lifting includes the identification of its reduction.
The actual-lifting torsor, and its triviality over a local noetherian
base, are recorded explicitly in the proof of
\cite[Lemma~24, p.590]{VasiuZink2010Purity}. Spread local liftings,
and use $H^1(\operatorname{Spec}R_0,\mathcal F)=0$ to choose an actual
lifting $B/R$. There are no infinitesimal automorphisms of an abelian
scheme as a group scheme, so this torsor of isomorphism classes is
an honest Zariski sheaf and its gluing is effective.

For each $r\geq1$, apply Serre--Tate to
\[
 R/(J+p^rR)\longleftarrow R/p^rR.
\]
It gives the unique lifting $A_r$ having BT group $M\bmod p^r$ and
the specified reduction of $A_0$. Relative to $B\bmod p^r$, its
difference is an element
\[
 d_r\in\Gamma\!\left(\operatorname{Spec}R_0,
   \operatorname{Lie}(A_0)\otimes\operatorname{Lie}(A_0^\vee)
        \otimes J/(J\cap p^rR)\right).
\]
These elements are compatible. The induced filtration
$J\cap p^rR$ is cofinal with the $p$-adic filtration on $J$ by
Artin--Rees. Since $J$ and the displayed coherent module are complete,
$(d_r)_r$ comes from a unique $d\in\Gamma(\mathcal F)$.
Translate the actual lifting $B$ by $d$. The resulting actual abelian
scheme $A/R$ has all the prescribed reductions $A_r$.

For each fixed $m$, the compatible isomorphisms
$A[p^m]\bmod p^r\simeq M[p^m]\bmod p^r$ lift uniquely to an
isomorphism over $R$: their coordinate algebras are finite projective
$R$-modules and hence complete. The Hopf identities, inverses and
compatibilities in $m$ are finite equalities and pass to the limit.
This gives $A[p^\infty]\simeq M$. If two liftings have the same data,
their difference in $\Gamma(\mathcal F)$ is zero modulo every
$J\cap p^rR$, hence is zero. Their reduction-preserving isomorphism
is unique by unramifiedness of the abelian Hom functor.

Finally induct along the finite filtration by powers of $J$.
Every step is a square-zero lifting of an actual abelian scheme,
which proves existence and uniqueness as stated.
\end{proof}

\begin{theorem}[C-RD12: exact nilpotent correction to the extension assertion]
\label{thm:C-RD12}
\status{Candidate proof; independent audit pending.}
Let $S$ be locally noetherian, let $T\subset V(p)$ be closed, and put
$U=S\setminus T$. Let $S_0\hookrightarrow S$ be a nilpotent closed
immersion. Fix a triple $(A_U,M,\varphi)$ as in C.1 and a specified
realization $A_0/S_0$ of its reduction.

Then this triple has a realization $A/S$ inducing the specified
$A_0$ if and only if the following condition holds on an affine open
cover of $S$. For $V=\operatorname{Spec}R$, choose a finitely generated
ideal $I$ containing $p$ with $T\cap V=V(I)$, and set
\[
 \widehat V=\operatorname{Spec}\widehat R_I,
 \qquad \widehat U=\widehat V\times_S U.
\]
Let $\widehat A/\widehat V$ be the unique algebraic lifting of
$A_0\times_{S_0}(\widehat V\times_SS_0)$ with BT group
$M_{\widehat V}$ supplied by Lemma~\ref{lem:C-RD12-complete}.
There must be an isomorphism
\[
 \widehat A|_{\widehat U}\ \simeq\ A_U|_{\widehat U}
                                                        \tag{C-RD12.1}
\]
which induces the specified reduction and the specified BT
identification $\varphi$. Thus, once the reduced triple is realized,
the punctured-completion comparison is an exact condition, with no
additional polarization, excellence, or formal-algebraization
hypothesis. It is an additional condition: it fails for
Proposition~\ref{prop:C-RD5}.
\end{theorem}
\begin{proof}
The completion $\widehat R_I$ is noetherian and $p$-adically complete.
Indeed a $p$-adically Cauchy sequence is $I$-adically Cauchy, and the
ideals $p^r\widehat R_I$ are $I$-adically closed; separatedness follows
from the Krull intersection theorem, since $p\in I$ lies in the
Jacobson radical. The kernel of
$\widehat R_I\to\widehat R_I/J\widehat R_I$, where $J$ defines
$S_0$, is nilpotent. Lemma~\ref{lem:C-RD12-complete} therefore applies
to this kernel; it is not necessary that it equal the entire
nilradical of $\widehat R_I$.

Necessity is base change and the uniqueness in that lemma. For
sufficiency, first suppose $J^2=0$ and work on $V=\operatorname{Spec}R$.
As in the lemma, choose an actual lifting $B/V$ of $A_0|_{V\cap S_0}$.
Let $\mathcal F$ be its coherent difference module. The schemes
$A_U|_{U\cap V}$ and $\widehat A$ have respective differences
\[
 d_U\in\Gamma(U\cap V,\mathcal F),\qquad
 d_{\widehat V}\in\Gamma(\widehat V,\mathcal F_{\widehat V})
\]
from the corresponding restrictions of $B$. Condition (C-RD12.1)
says exactly that these sections agree on $\widehat U$.
The quasi-coherent gluing equivalence of Ferrand--Raynaud, in the
precise form \cite[Theorem~0.3, p.561]{MoretBailly1996Descent}, gives
\[
 \Gamma(V,\mathcal F)\ \simeq\
 \Gamma(U\cap V,\mathcal F)
 \mathop{\times}_{\Gamma(\widehat U,\mathcal F)}
 \Gamma(\widehat V,\mathcal F_{\widehat V}).
\]
Translate $B$ by the resulting section. This produces an actual
abelian scheme $A/V$ with the required isomorphisms over $U\cap V$
and $\widehat V$.

Its BT group has the required comparison with $M$. At each finite
level, descend the compatible Hopf-algebra isomorphisms over
$U\cap V$ and $\widehat V$ by the same quasi-coherent equivalence.
Their group laws and compatibility between levels descend as
identities. Thus $A[p^\infty]\simeq M$, with the prescribed
identification over $U\cap V$.

The construction is unique among liftings with the specified
reduction and these restrictions: its difference from another is a
section of $\mathcal F$ vanishing on $U\cap V$ and $\widehat V$,
hence zero by the displayed equalizer. On overlaps the same argument,
after restriction to smaller affine opens, provides the unique
reduction-preserving isomorphisms, so the schemes glue. Induction
along the finite powers of $J$ gives the general nilpotent case.
At each stage the two prescribed restrictions are the reductions of
$A_U$ and $\widehat A$, so (C-RD12.1) provides the required compatible
differences. This proves sufficiency globally.
\end{proof}

\begin{remark}[What this criterion does and does not remove]
Taking $S_0=S_{\mathrm{red}}$ separates the essential-image problem
into the reduced realization problem and an exact additional
nilpotent comparison. Theorem~\ref{thm:C-RD11} shows that the first
problem can already fail under a) and b). Consequently one cannot
replace it by the assumption ``$S$ is reduced''. The theorem above
settles the extra condition for nonreduced bases once a reduced
realization is specified; it does not assert automatic effectivity
of a general unpolarized formal abelian scheme.
\end{remark}
